\chapter{启动、链接脚本与内核初始化}

\section{原理}

从裸机程序走向内核，第一道门是启动。机器上电后不会自动理解 C 语言，也不会自动准备堆、栈、全局变量和设备。CPU 只知道从复位向量或固件指定入口取指。操作系统必须把“二进制文件如何进入内存、第一条指令在哪里、栈在哪里、全局变量如何初始化、设备何时可用”这些问题全部讲清。

在单片机上，启动路径通常是 reset handler：设置栈指针，把 \code{.data} 从 Flash 复制到 SRAM，把 \code{.bss} 清零，初始化时钟和最小串口，然后调用 \code{main} 或 \code{kernel\_main}。在 PC 或服务器上，路径更长：固件初始化硬件，bootloader 加载内核镜像和启动参数，切换 CPU 模式，建立最小页表，最后跳到内核入口。

启动阶段的特点是“没有服务可用”。还没有调度器，不能睡眠；还没有内存分配器，不能随便 \code{malloc}；还没有完整中断，不能依赖异步事件；还没有文件系统，不能读普通文件；还没有日志系统，调试输出只能靠串口、早期 console 或内存缓冲区。这些约束解释了为什么内核初始化代码通常按严格阶段组织。

\subsection{上电后 CPU 为什么不能直接运行内核 C 函数}

C 函数依赖一组运行环境：栈指针有效，全局变量已初始化，BSS 为零，代码和数据位于链接器预期地址，调用约定可用，必要的地址映射存在。上电瞬间这些条件都不自动成立。CPU 只处在架构定义的复位状态，从固定入口取指；剩下的环境必须由启动代码一步步建立。

\begin{longtable}{p{0.22\linewidth}p{0.34\linewidth}p{0.32\linewidth}}
\toprule
C 代码假设 & 谁建立 & 若缺失会怎样 \\
\midrule
栈可用 & reset handler 或 bootloader & 函数调用和局部变量破坏未知内存 \\
全局变量已初始化 & 启动代码复制 \code{.data}、清 \code{.bss} & 锁、队列、指针从随机值开始 \\
代码地址正确 & 链接脚本和装载器 & 跳转、全局地址、异常表错误 \\
内存可访问 & 页表或物理映射 & 取指或访存 fault \\
最小输出可用 & early console 初始化 & 出错时没有证据 \\
\bottomrule
\end{longtable}

所以内核入口往往先是汇编或极小 C 代码。它不是为了炫技，而是因为在运行时环境建立前，普通 C 抽象还没有可靠落脚点。

\subsection{链接脚本决定内存布局}

链接脚本把代码段、只读数据、已初始化数据、未初始化数据和栈安排到具体地址。没有链接脚本，编译器和链接器不知道内核应该放在物理地址哪里，也不知道哪些符号表示段边界。

\begin{lstlisting}
SECTIONS {
  . = 0x80200000;
  __kernel_start = .;
  .text : { *(.text*) }
  .rodata : { *(.rodata*) }
  .data : { *(.data*) }
  __bss_start = .;
  .bss : { *(.bss*) *(COMMON) }
  __bss_end = .;
  __kernel_end = .;
}
\end{lstlisting}

这些符号不是装饰。启动代码需要 \code{\_\_bss\_start} 和 \code{\_\_bss\_end} 来清零 BSS；早期物理页分配器需要 \code{\_\_kernel\_end} 知道哪些内存已经被内核占用；页表初始化要知道哪些区域可执行、只读或可写。

\subsection{链接地址、加载地址和重定位}

链接地址是代码认为自己运行的位置，加载地址是 bootloader 把镜像放到内存的位置。二者可以相同，也可以不同。若不同，早期代码要么以位置无关方式运行，要么完成重定位，要么建立页表把链接虚拟地址映射到实际物理地址。

\begin{lstlisting}
kernel linked at virtual address: 0xffffffff81000000
kernel loaded at physical address: 0x00100000

early page table maps:
  0xffffffff81000000 -> 0x00100000
\end{lstlisting}

这解释了为什么 64 位内核常有“高半区内核”概念。用户进程使用低地址范围，内核映射在高虚拟地址范围；每个进程切换页表时仍保留内核映射，系统调用和中断进入内核后能访问同一套内核代码和数据。这个设计依赖页表，而不是物理内存真的在高地址。

\subsection{启动阶段的状态机}

内核启动可以写成状态机：

\begin{lstlisting}
Reset
  -> MinimalStack
  -> DataCopied
  -> BssCleared
  -> EarlyConsole
  -> BootMemoryDiscovered
  -> PageTablesReady
  -> InterruptsReady
  -> SchedulerReady
  -> InitProcessStarted
\end{lstlisting}

每个阶段只能使用之前已经建立的能力。比如 EarlyConsole 之前不能输出复杂日志；BootMemoryDiscovered 之前不能建立通用页分配；SchedulerReady 之前不能阻塞等待；InterruptsReady 之前不能依赖设备中断。

\subsection{x86-64 模式演进：实模式、保护模式、长模式}

x86-64 的历史兼容层让启动路径明显分阶段。复位后 CPU 从兼容模式开始，早期代码或固件/bootloader 需要逐步建立 GDT、打开保护机制、准备页表、启用 PAE 和长模式，最后跳到 64 位代码。现代 UEFI 和 bootloader 会屏蔽许多细节，但内核仍要理解当前进入时的模式和页表状态。

\begin{lstlisting}
legacy shape:
  reset vector
  real mode firmware/bootloader
  load kernel setup code
  build provisional GDT
  enter protected mode
  build early page tables
  enable long mode
  jump to 64-bit kernel entry
\end{lstlisting}

GDT 在长模式下不再像 32 位分段那样用于普通地址计算，但仍参与代码段、数据段、TSS、特权切换和系统调用入口设置。TSS 也不只是历史名词，x86-64 可用 IST 为 NMI、double fault 等关键异常提供专用栈，避免在当前栈损坏时继续崩溃。

\subsection{BIOS、UEFI 和 bootloader 的职责边界}

BIOS 风格启动和 UEFI 风格启动不同，但目标相同：把内核镜像和硬件描述交给内核。bootloader 不应替内核做所有工作，它只建立足够条件让内核接管。内核接管后，必须停止依赖大部分固件 Boot Services，改用自己的内存管理、驱动和中断。

\begin{longtable}{p{0.2\linewidth}p{0.34\linewidth}p{0.34\linewidth}}
\toprule
阶段 & 可依赖对象 & 退出后的边界 \\
\midrule
固件早期 & DRAM 初始化、平台基础设置 & 内核通常看不到细节，只接收结果 \\
bootloader & 读取磁盘/网络、加载镜像、构造参数 & 不再参与普通运行时服务 \\
UEFI Boot Services & 分配内存、访问协议、获取内存地图 & \code{ExitBootServices} 后不可再调用 \\
内核 early boot & 启动参数、内存地图、早期 console & 必须校验并建立自己的对象 \\
\bottomrule
\end{longtable}

\code{ExitBootServices} 是一个硬边界。退出前，固件仍可能拥有内存和设备；退出后，内核承担管理责任。若内核拿到内存地图后，固件又分配了内存，而内核没有重新获取地图，就可能把固件仍使用的页当成空闲页。

\section{实践}

实践时先写一个启动依赖表：

\begin{longtable}{p{0.24\linewidth}p{0.32\linewidth}p{0.32\linewidth}}
\toprule
阶段 & 可用能力 & 禁止事项 \\
\midrule
Reset & 寄存器、固定入口 & 使用全局变量假设已初始化 \\
MinimalStack & 栈、少量汇编 & 调用复杂 C 运行库 \\
BssCleared & 零初始化全局变量 & 动态分配内存 \\
EarlyConsole & 最小输出 & 格式化复杂日志、大缓冲 \\
BootMemory & 物理内存范围 & 使用完整虚拟内存假设 \\
PageTables & 基础地址翻译 & 访问未映射设备 \\
Scheduler & 任务切换 & 在早期 init 中遗留永久锁 \\
\bottomrule
\end{longtable}

然后画出内核镜像布局：text、rodata、data、bss、boot stack、page tables、free memory。只要能解释“内核自己占了哪些页，哪些页可以交给分配器”，后续物理页管理才有根。

\section{算法与实现分析}

\subsection{BSS 清零和 data 复制}

\begin{lstlisting}
boot_zero_bss():
  p = __bss_start
  while p < __bss_end:
    *p = 0
    p += word_size

boot_copy_data():
  src = __data_load_start
  dst = __data_start
  while dst < __data_end:
    *dst++ = *src++
\end{lstlisting}

这两个循环是 \code{O(n)}，n 是段大小。它们必须在使用对应全局变量之前完成。若 BSS 没清零，全局队列、锁和指针会从随机值开始，后续 bug 很难定位。

\subsection{早期 bump allocator}

完整页分配器建立前，内核常用 bump allocator 分配早期对象。它只维护一个 next 指针，每次按对齐向上推进，不支持释放。

\begin{lstlisting}
early_alloc(size, align):
  next = align_up(early_next, align)
  end = next + size
  if end > early_limit: panic()
  early_next = end
  return next
\end{lstlisting}

分配是 \code{O(1)}，实现极简单。代价是不能释放，且必须在切换到正式分配器后停止使用。它适合页表、启动参数副本、早期设备表等生命周期贯穿启动阶段的对象。

\subsection{初始化依赖图}

内核初始化不能只靠函数排列顺序。一个模块可能依赖 console，一个模块依赖内存分配器，一个模块依赖中断，一个模块必须在调度器启动前完成。把这些关系写成依赖图，才能解释为什么初始化常被分成 early、core、subsys、device、late 等阶段。

\begin{lstlisting}
init_node {
  name: "page_allocator"
  requires: ["boot_memory_map", "early_console"]
  provides: ["alloc_pages"]
}

run_init_graph(nodes):
  ready = nodes with all requirements satisfied
  while ready not empty:
    n = pop(ready)
    run(n.init)
    mark_provided(n.provides)
    add newly unblocked nodes to ready
  if unfinished nodes remain:
    panic("init dependency cycle or missing provider")
\end{lstlisting}

这个算法是拓扑排序，复杂度是 \code{O(V + E)}。它的工程价值是把“某函数必须先跑”变成可检查约束。若出现环，例如文件系统初始化依赖块设备，而块设备初始化又要从文件系统读取配置，就必须拆出更早的配置来源。

\subsection{启动参数与可信边界}

bootloader 会把 UEFI 内存地图、命令行、initrd 和 ACPI 表交给内核。内核不能盲信这些数据，因为它们可能来自固件 bug、配置错误或攻击面。早期解析必须检查地址范围、长度、对齐和重叠。

\begin{lstlisting}
validate_memory_region(region):
  require region.start aligned to PAGE_SIZE
  require region.length > 0
  require checked_add(region.start, region.length, &end)
  require end <= physical_address_limit
  require not overlaps(kernel_image)
  require not overlaps(previous_accepted_regions)
\end{lstlisting}

若把内核镜像所在物理页误标为空闲，后续页分配器会把正在运行的内核代码分配给别人。启动阶段的错误通常不是立刻崩溃，而是在很久以后变成随机损坏，因此必须在早期就把内存地图打印和校验做好。

\subsection{ACPI 和 PCI 配置空间：硬件不是凭空出现的}

内核要知道“有哪些 CPU、哪些内存、哪些设备、设备寄存器在哪、用哪个 IRQ、挂在哪个 NUMA node、有哪些电源关系”。在本书的 x86-64 主线里，这部分信息主要来自 UEFI/ACPI、PCI/PCIe 配置空间和固件内存地图。它们都不是设备驱动本身，而是硬件描述入口。

\begin{longtable}{p{0.2\linewidth}p{0.34\linewidth}p{0.34\linewidth}}
\toprule
描述来源 & x86-64 常见用途 & 描述内容 \\
\midrule
ACPI & PC、服务器 & CPU 拓扑、NUMA、PCI routing、电源、热管理、平台设备资源 \\
PCI config & 可枚举 PCIe 设备 & vendor/device、BAR、MSI、能力链 \\
UEFI memory map & 启动阶段内存交接 & RAM、保留区、runtime services、ACPI reclaim 区域 \\
\bottomrule
\end{longtable}

驱动 probe 的输入很多来自这些描述。若 ACPI resource 或 PCI BAR 错，MMIO 映射到错误地址；若 IRQ routing 或 MSI-X 配置错，设备完成不会唤醒驱动；若 DMA mask 或 IOMMU domain 配置错，设备可能访问不到高地址内存。硬件描述错误会表现成软件 bug，所以启动阶段要尽早暴露和校验资源表。

\subsection{从单核启动到多核启动}

多核系统通常只有 bootstrap processor 先运行内核入口，其他 CPU 处于等待状态。主 CPU 建立全局页表、per-CPU 区域和调度数据后，再唤醒 application processors。每个 CPU 都需要自己的内核栈、当前任务指针、runqueue 和中断状态。

\begin{lstlisting}
boot_cpu_main():
  setup_global_memory()
  setup_percpu_area(boot_cpu)
  for cpu in secondary_cpus:
    allocate_cpu_stack(cpu)
    send_startup_ipi(cpu, secondary_entry)
  wait_until_all_online()
  start_scheduler()

secondary_entry(cpu):
  load_kernel_page_table()
  set_percpu_base(cpu)
  init_local_interrupts()
  mark_cpu_online(cpu)
  idle_loop()
\end{lstlisting}

这里的不变量是：CPU online 之前不能被调度器选作迁移目标；per-CPU 数据初始化之前不能接收普通设备中断；所有 CPU 使用的内核文本映射必须一致。多核启动把“初始化顺序”从单线流程变成了并发协议。

\subsection{per-CPU 数据为什么必须很早准备}

多核 OS 不能让所有 CPU 频繁写同一批全局变量。每个 CPU 需要自己的当前任务指针、内核栈、runqueue、计数器、本地 timer 状态、抢占计数和中断统计。这些对象通常放在 per-CPU 区域，通过特殊基址寄存器或固定偏移快速访问。

\begin{lstlisting}
per_cpu_area[cpu]:
  current_task
  kernel_stack_top
  runqueue
  irq_count
  preempt_count
  local_timer_state
\end{lstlisting}

secondary CPU 进入内核后，必须先设置自己的 per-CPU base，才能安全访问 \code{current} 这类隐式变量。否则 CPU1 可能误用 CPU0 的当前任务或 runqueue，造成调度和中断统计混乱。per-CPU 是多核性能优化，也是正确性边界。

\section{测试}

\begin{itemize}
\tightlist
\item 链接脚本导出的段边界单调递增且按页对齐。
\item BSS 清零后，全局计数器、指针和队列初值可预测。
\item 早期分配器返回地址满足对齐，且不会覆盖内核镜像。
\item 未初始化 console 前，日志路径不会访问空设备。
\item 调度器 ready 前，任何初始化函数都不能睡眠。
\item 页表 ready 后，text 只读可执行，data/bss 可写不可执行。
\end{itemize}

启动章节的验收标准是：能从第一条指令解释到第一个用户任务启动，中间每一阶段可用能力和禁止事项都清楚。

\section{源码级补充：x86-64、UEFI 与 start\_kernel}

\subsection{RESET 到长模式}

x86 CPU 复位后不是直接处于 64 位长模式。传统路径会经历实模式、保护模式和长模式切换；现代 UEFI 固件和 bootloader 会替内核完成大量工作，但内核仍必须理解启动时 CPU 模式、页表、内存地图和调用约定。

\begin{lstlisting}
RESET
  -> firmware initializes platform
  -> bootloader loads kernel image and initrd
  -> build boot parameters and memory map
  -> enter protected/long mode as required by protocol
  -> jump to decompressor or kernel entry
  -> early page tables
  -> start_kernel()
\end{lstlisting}

实模式只能直接访问有限地址空间，保护模式引入段和保护，长模式引入 64 位寄存器和分页模型。操作系统教材不需要把每条切换指令都背下来，但必须知道：64 位内核能运行，是因为早期代码已经设置好 GDT、控制寄存器、EFER.LME、页表和跳转序列。

\subsection{BIOS 与 UEFI 的差异}

BIOS 风格启动常从磁盘引导扇区开始，bootloader 逐步加载更多代码。UEFI 则提供标准化的 Boot Services、Runtime Services、内存地图和 PE/COFF 应用加载环境。UEFI 下 bootloader 或 EFI stub 能以更结构化方式获得 GOP framebuffer、ACPI 表、initrd 和命令行。

\begin{longtable}{p{0.2\linewidth}p{0.34\linewidth}p{0.34\linewidth}}
\toprule
项目 & BIOS 风格 & UEFI 风格 \\
\midrule
入口形态 & 引导扇区和 bootloader 阶段 & EFI 应用或 EFI stub \\
固件服务 & 中断调用，模式受限 & Boot Services/Runtime Services \\
内存地图 & e820 等接口 & \code{GetMemoryMap} \\
图形输出 & VGA/vesa 等路径 & GOP framebuffer \\
退出固件 & bootloader 控制移交 & \code{ExitBootServices} 后服务受限 \\
\bottomrule
\end{longtable}

\code{ExitBootServices} 是重要边界。退出前可以调用 UEFI Boot Services 分配内存、打开协议、读取设备；退出后内核必须依靠自己的驱动和内存管理。若退出前后内存地图不同步，内核可能误用固件保留页。

\subsection{Linux start\_kernel 调用序列}

Linux 的 \code{start\_kernel} 是阅读大型内核初始化的主干。它不是所有细节的开始，但它把架构初始化、内存、调度、RCU、中断、定时器、VFS、proc、driver model、init 进程串起来。

\begin{lstlisting}
start_kernel()
  setup_arch()
  setup_command_line()
  mm_init()
  sched_init()
  radix_tree_init()
  trap_init()
  init_IRQ()
  tick_init()
  rcu_init()
  vfs_caches_init()
  rest_init()
\end{lstlisting}

不同版本顺序会变化，读源码时应抓依赖：\code{setup\_arch} 解析固件和内存；\code{mm\_init} 建立内存分配基础；\code{sched\_init} 建立 runqueue；\code{trap\_init/init\_IRQ} 建立入口和中断；\code{rest\_init} 创建 kernel thread 和第一个用户态 init 的路径。

\subsection{initcall 机制}

驱动和子系统不能全都手写进一个巨大初始化函数。Linux 用 initcall 等级把初始化函数放入特殊段，启动时按等级依次调用。常见等级包括 early、core、postcore、arch、subsys、fs、device、late。

\begin{lstlisting}
early_initcall(foo_early);
core_initcall(mm_subsystem);
subsys_initcall(bus_subsystem);
fs_initcall(vfs_feature);
device_initcall(driver_init);
late_initcall(debug_late);
\end{lstlisting}

initcall 的实质是链接脚本和函数指针数组。它让模块能声明“我属于哪个初始化阶段”，但不能自动解决所有依赖。若一个 device initcall 假设某总线已注册，而总线初始化等级写错，启动会出现只在特定配置下复现的失败。

\subsection{多核启动和 CPU hotplug}

bootstrap processor 先运行全局初始化，然后通过 IPI 或固件接口唤醒 application processors。每个 CPU 上线前要设置 per-CPU base、内核栈、本地 APIC、idle task 和 runqueue。CPU hotplug 则要求运行中 CPU 可下线和上线，所以初始化/清理路径必须可重复、可回滚。

\begin{lstlisting}
bringup_cpu(cpu):
  allocate_idle_task(cpu)
  allocate_percpu_area(cpu)
  send_startup_ipi(cpu)
  wait cpu reports online
  add cpu to scheduler domains
  enable interrupts on cpu
\end{lstlisting}

失败案例：CPU 已经标记 online 但本地 timer 未初始化，调度器可能把任务迁过去后无法抢占；CPU 下线时未迁走 timer 或 workqueue，后续事件投递到离线 CPU。多核启动不是“让其他核心开始跑”，而是把每个 CPU 纳入调度、中断和时间子系统的一致状态。
